\section{Theorems and proofs for the partial results}\label{app:partial}

\subsection{Problem 16.14}\label{app:1614}
Notation: $G$ is a finite 2-group with $\Omega_1(G)\le Z(G)$; $E=\Omega_1(G)$ (elementary abelian, central,
generated by all involutions), $r=\rk E=\rk Z(G)$, $d=d(G)=\rk(G/G^2)=\dim G/\Phi(G)$.

\begin{theorem}[class $\le 2$]\label{thm:1614A}
If $G$ has class $\le 2$ and $G^2\le Z(G)$, then $d(G)\le 2\,\rk Z(G)$.
\end{theorem}
\begin{proof}
Put $Z=Z(G)$ and $r=\rk Z=\dim(Z/Z^2)$. Since $G$ has class $\le 2$, $(gh)^2=g^2h^2[h,g]$ with $[h,g]\in G'\le Z$.
Hence $Q:G\to Z/Z^2$, $g\mapsto g^2Z^2$, satisfies $Q(gh)=Q(g)Q(h)[h,g]Z^2$, $Q(gz)=Q(g)$ for $z\in Z$, $Q(g^2)=1$.
So $Q$ factors through the $\F_2$-space $V=G/(ZG^2)$ and is a quadratic map $q:V\to W=Z/Z^2\cong\F_2^r$ with polar
form the commutator pairing; in coordinates $q=(q_1,\dots,q_r)$. If $v\in V\setminus\{0\}$ had $q(v)=0$, a
preimage $g\notin ZG^2$ would satisfy $g^2=z^2$ with $z\in Z$, so $(gz^{-1})^2=1$, $gz^{-1}\in\Omega_1(G)\le Z$ and
$g\in Z$, a contradiction. So $q_1,\dots,q_r$ have no common non-trivial zero, and by the Chevalley--Warning theorem
$n:=\dim V\le 2r$. With $W'=\operatorname{span}q(V)$ and $r'=\dim W'\le r$ the same argument gives $n\le 2r'$.
Finally $\Phi(G)=G^2G'\le Z$, $Z^2\le G^2\le\Phi(G)$ and $\Phi(G)Z^2/Z^2\supseteq W'$, so
$d(G)=\dim G/ZG^2+\dim ZG^2/\Phi(G)=n+\dim Z/\Phi(G)\le n+(r-r')\le 2r'+r-r'\le 2r$.
\end{proof}
If $\Phi(G)$ is central and elementary abelian ($p$-class $\le 2$), the theorem applies.

\begin{theorem}[minimal counterexample]\label{thm:1614B}
Let $G$ be a counterexample of minimal order ($d(G)>2r(G)$ and $d(K)\le 2r(K)$ for all smaller $K$ with
$\Omega_1(K)\le Z(K)$), $E=\Omega_1(G)$, $r=r(G)$, $H=G/E$. Then:
\begin{enumerate}[label=(\roman*)]
\item $d(G)=2r+1$, and every maximal subgroup $M$ satisfies $\Omega_1(M)=E$ and $d(M)=2r=2r(M)$;
\item $E\le\Phi(G)$, hence $d(H)=2r+1$;
\item $\Phi(M)=\Phi(G)$ for every maximal $M$, hence $\Phi(\bar M)=\Phi(H)$ for every maximal $\bar M\le H$;
\item every $g\in G\setminus E$ has order $2\cdot|gE|$;
\item for every elementary abelian $A\le H$ the squaring map induces a quadratic map $q_A:A\to E$ without
non-trivial zeros; hence $\rk A\le 2r$;
\item every subgroup of $H$ with trivial Schur multiplier is abelian;
\item $r\ge 2$, $|\Phi(G)|\ge 2|E|$ and $|G|\ge 2^{3r+2}$.
\end{enumerate}
\end{theorem}
\begin{proof}
Subgroups inherit the hypothesis: for $K\le G$, $\Omega_1(K)\le E\cap K\le Z(K)$ and $r(K)\le r$. For a maximal
subgroup $M$, $M/\Phi(G)$ is an elementary abelian quotient of $M$ of rank $d-1$, so $d(M)\ge d(G)-1$; by
minimality $d(M)\le 2r(M)\le 2r$, hence $d(G)\le 2r+1$, so $d(G)=2r+1$, $d(M)=2r$, $r(M)=r$ and $\Omega_1(M)=E$: (i).
(ii) If $t\in E\setminus\Phi(G)$, pick a maximal $M$ with $t\notin M$; then $G=M\times\gen t$, $r(M)=r-1$ and
$d(M)=2r=2r(M)+2$, contradicting minimality for $M$. (iii) $\Phi(M)\le\Phi(G)$ ($M$ normal) and
$|M:\Phi(M)|=2^{2r}=|M:\Phi(G)|$. (iv) Let $|gE|=2^k$; then $g^{2^k}\in E$ and $g^{2^{k-1}}\notin E$; if $g^{2^k}=1$
then $g^{2^{k-1}}$ would be an involution outside $E$. (v) The preimage $\tilde A$ of $A$ has $\tilde A^2\le E$,
$\tilde A'\le E\le Z(\tilde A)$, so $q_A(aE)=a^2$ is a well-defined quadratic map with polar form the commutator
pairing; $q_A(aE)=0$ with $aE\ne 0$ would give an involution $a\notin E$. Chevalley--Warning gives $\rk A\le 2r$.
(vi) For $Q\le H$ with preimage $\tilde Q$, the five-term sequence $H_2(\tilde Q)\to H_2(Q)\to E\to\tilde Q^{ab}\to
Q^{ab}\to 0$ gives $E\cap\tilde Q'=\mathrm{im}(H_2(Q)\to E)=0$, so $\tilde Q'\cong Q'$; if $Q'\ne 1$ the 2-group
$\tilde Q'$ contains an involution outside $E$. (vii) $r=1$ would make $G$ cyclic or generalised quaternion
($d\le 2$); $\Phi(G)=E$ would make $\Phi(G)$ central elementary abelian and Theorem~\ref{thm:1614A} would apply.
So $|\Phi(G)|\ge 2^{r+1}$ and $|G|=2^d|\Phi(G)|\ge 2^{3r+2}$.
\end{proof}

\begin{theorem}[small orders]\label{thm:1614C}
There is no counterexample of order $\le 2^{11}$, and none of order $2^{12}$ with $r\ge 3$.
\end{theorem}
\begin{proof}[Proof (computation guided by Theorem~\ref{thm:1614B})]
By (vii) a minimal counterexample has $|G|=2^{2r+1}|\Phi(G)|$ with $|\Phi(G)|\ge 2^{r+1}$, $r\ge 2$: for $r=2$,
$|G|\ge 2^8$; for $r=3$, $|G|\ge 2^{11}$; for $r\ge 4$, $|G|\ge 2^{14}$. Order $2^8$ is excluded by the direct sweep of
all groups of order 256. For the remaining cases $H=G/E$ is a group of order $|G|/2^r\in\{2^7,2^8,2^9\}$ with
$d(H)=2r+1$, available in the SmallGroups library, and $G$ is a central extension of $H$ by $E=\F_2^r$ with
$E\le\Phi(G)$ in which, by (iv), no involution of $H$ lifts to an involution. Every such $G$ is a quotient of the
2-covering group $P^*=F/R^2[F,R]$ of $H=F/R$. Hence $H$ is excluded as soon as some involution of $H$ lifts to an
involution of $P^*$ (filter F3$'$, computed with ANUPQ's \code{PqPCover}; it is equivalent to the vanishing on all
2-cocycles of the functional ``square of the lift'', which was used in the first version of the computation), or
some maximal subgroup $\bar M$ of $H$ has $\Phi(\bar M)\ne\Phi(H)$ (filter F2, by (iii)), or $H$ has an elementary
abelian subgroup of rank $2r+1$ (filter F1, by (v)). Surviving $H$ would require a search for $r$ cohomology classes
satisfying the lifting condition; none survived. The cases: $r=3$, $|G|=2^{11}$: the 10 groups of order 256 with
$d=7$ (two fail F1, eight F3); $r=3$, $|G|=2^{12}$: the 1139 groups of order 512 with $d=7$ (32 fail F1, 82 F2,
1025 F3); $r=2$, $|G|=2^9$: the 169 groups of order 128 with $d=5$; $r=2$, $|G|=2^{10}$: the 28\,653 groups of order 256
with $d=5$ (7043 fail F2, 21\,610 F3); $r=2$, $|G|=2^{11}$: all $7\,111\,878$ groups of order 512 with $d=5$ (1049 fail
F2, $7\,110\,829$ F3$'$). Logs: \code{work/p16\_14/r256log*}, \code{r512log\_*}, \code{r128log}, \code{r256r2log\_*},
\code{r512r5log\_*}.
\end{proof}

\subsection{Problem 20.37}\label{app:2037}
Let $L(G)=\{|A|:G=AB\text{ for some }A,B\subseteq G\}$. (a) $G=AB$ with $|A||B|=|G|$ iff all products $ab$ are
distinct iff $A^{-1}A\cap BB^{-1}=\{1\}$; (b) $L(G)$ is closed under $a\mapsto|G|/a$ (as $G=B^{-1}A^{-1}$); (c) for
$H\le G$ with left transversal $T$ and right transversal $T'$, $G=TH=HT'$, so $L(G)\supseteq L(H)\cup|G:H|L(H)$. The
chain-achievable set $L'(G)$ is the closure of $\{1,|G|\}$ under (b) and (c) over all subgroups (maximal subgroups
recursively suffice).

\begin{lemma}\label{lem:2037L1}
Let $H\le G$, $|H|=h$, and $x_1=1,x_2,\dots,x_k\in G$ with distinct right cosets $Hx_i$; put $A=\bigcup Hx_i$. If
$B\subseteq G$ is such that the $k$-sets $\{Hx_1b,\dots,Hx_kb\}$ ($b\in B$) partition $H\backslash G$, then $G=AB$.
\end{lemma}
\begin{proof}
$AB=\bigcup_{i,b}Hx_ib$ is the union of all right cosets, each once, so $AB=G$ and $|A||B|=|G|$.
\end{proof}

\begin{theorem}\label{thm:2037T1}
Let $H\le K\le G$ with $|K:H|$ even and $K=\gen{H,x}$ for some $x\in K$ (necessarily $x\notin H$). Then there is
$B\subseteq K$ with $|B|=|K|/(2|H|)$ and $K=(H\cup Hx)B$; in particular $2|H|\in L(K)\subseteq L(G)$.
\end{theorem}
\begin{proof}
Let $\Gamma$ be the simple graph with vertex set $H\backslash K$ and an edge $\{Hb,Hxb\}$ for every $b\in K$ with
$Hb\ne Hxb$. $K$ acts on $H\backslash K$ by right multiplication, mapping the edge $\{Hb,Hxb\}$ to $\{Hbg,Hxbg\}$,
so $K$ acts on $\Gamma$ by automorphisms, transitively on vertices. The neighbours of $Hb$ are the cosets
$Hx^{\pm1}hb$ ($h\in H$), so the vertices reachable from $Hb$ are the $Hwb$ with $w\in\gen{H,x}=K$: $\Gamma$ is
connected. By the Gallai--Edmonds structure theorem~\cite{LP}, the set $D$ of vertices missed by some maximum
matching is $\Aut(\Gamma)$-invariant, hence $\emptyset$ or all of $V(\Gamma)$; in the latter case the connected
graph $\Gamma$ would be factor-critical and have an odd number of vertices, contradicting $|K:H|$ even. So $D=\emptyset$
and $\Gamma$ has a perfect matching $M$. Choosing $b_e$ with $e=\{Hb_e,Hxb_e\}$ for each $e\in M$ and $B=\{b_e\}$,
the 2-sets $\{Hb,Hxb\}$ ($b\in B$) partition $H\backslash K$, and Lemma~\ref{lem:2037L1} applied in $K$ gives
$K=(H\cup Hx)B$.
\end{proof}

\begin{lemma}[two-sided construction]\label{lem:2037L2}
Let $H,K\le G$ with all $(H,K)$-double cosets regular ($|HgK|=|H||K|$ for all $g$), $\Omega=H\backslash G/K$,
$X=\{x_1,\dots,x_k\}$ and $Y=\{y_1,\dots,y_l\}\subseteq G$ such that the $k$-sets $\{Hx_iy_jK:i\}$ ($j=1,\dots,l$)
partition $\Omega$. Then $G=AB$ with $A=HX$ and $B=YK$, $|A|=k|H|$, $|B|=l|K|$.
\end{lemma}
\begin{proof}
$AB=\bigcup_{i,j}Hx_iy_jK$ is the union of all double cosets, each once; by regularity $|AB|=kl|H||K|=|A||B|$.
\end{proof}
In the computations, $X$ is chosen at random and $Y$ is found by an exact-cover search on $\Omega$ (a dancing-links
solver with node limits, repeated with new $X$); all found pairs $(A,B)$ are verified by direct multiplication in
GAP (\code{work/p20\_37/verify\_all2.g}: 286 factorisations, 0 failures).

\subsection{Problem 20.21}\label{app:2021}
Let $G$, $K\cong L$ be as in the problem and $I=K\cap L$. Then $G/I$ embeds in $C_{12}\times A_4$, $|K/I|=|L/I|\in
\{4,12\}$, and the case 12 leads, via an isomorphism $\varphi:L\to K$, to a smaller configuration of the same kind
(so an example of minimal order is not of this type, and every example contains one of the case-4 type). In the case 4, $K/I\cong V_4$, $L/I\cong C_4$, $G/I\cong C_4\times A_4$,
and an element $s$ mapping to a 3-cycle acts on $K/I$ as a 3-cycle. With $I'=\varphi(I)\trianglelefteq K$ one has
$K/I'\cong C_4$ and $I'\cong I$. Hence any example satisfies the \emph{local condition} (L): a group $M$ with
normal subgroups $I,I'\cong I$, $M/I\cong V_4$, $M/I'\cong C_4$, and an automorphism $\alpha$ with $\alpha(I)=I$
permuting the three non-trivial cosets of $I$ cyclically.

\begin{lemma}\label{lem:2021ab}
No abelian group satisfies (L).
\end{lemma}
\begin{proof}
Induction on $|M|$. $\Omega_1(M)$ is characteristic, so its image in $M/I\cong V_4$ is $\alpha$-invariant, hence $1$
or $V_4$. As $M$ is abelian, $\Omega_1(M)\cap I=\Omega_1(I)$ and $\Omega_1(M)\cap I'=\Omega_1(I')$, so
$|\Omega_1(M)\cap I|=|\Omega_1(M)\cap I'|$. If $\Omega_1(M)I=M$ then $|\Omega_1(M)\cap I'|=|\Omega_1(M)|/4$, so
$\Omega_1(M)I'=M$ and $M/I'$ would be elementary abelian, contradicting $M/I'\cong C_4$. Hence
$\Omega_1(M)\le I\cap I'$, and $(M/\Omega_1(M);I/\Omega_1(I),I'/\Omega_1(I'),\bar\alpha)$ is a smaller configuration
(L) ($I/\Omega_1(I)\cong I'/\Omega_1(I')$ because $I\cong I'$).
\end{proof}

\begin{lemma}\label{lem:2021T2}
No configuration (L) has $II'\ne M$ (equivalently $|I:I\cap I'|=2$).
\end{lemma}
\begin{proof}
Induction on $|M|$; write $N_2(X)$ for the number of involutions in $X$. Since $\alpha$ permutes the cosets $aI,bI,cI$
cyclically, $N_2(aI)=N_2(bI)=N_2(cI)=:c_2$ and $N_2(M)=N_2(I)+3c_2$. With $\bar u$ generating $M/I'$, the cosets $uI'$,
$u^3I'$ contain no involutions, so $N_2(M)=N_2(I')+N_2(u^2I')$, and $N_2(I)=N_2(I')$ gives $N_2(u^2I')=3c_2$. Now $II'$
has index 2 in $M$; $II'/I'$ is the subgroup of order 2 of $M/I'$, so $II'=I'\cup u^2I'$, and $II'=I\cup aI$ for one
non-trivial coset. Counting involutions in $II'$ in two ways gives $N_2(I)+c_2=N_2(I')+3c_2$, so $c_2=0$: all
involutions lie in $I$ and in $I'$. Then $W=\Omega_1(M)=\Omega_1(I)=\Omega_1(I')\le I\cap I'$, $W\ne 1$, and
$(M/W;I/W,I'/W,\bar\alpha)$ is a smaller configuration (L) with $|I/W:(I\cap I')/W|=2$.
\end{proof}
Consequently an example of minimal order must have $K=I\cdot\varphi(I)$. Computationally (\code{work/p20\_21}): no
group of order $\le 2000$ in an order divisible by 12 with $\le 30\,000$ groups has the required normal subgroups;
and no group of order $\le 256$, of order 512 and rank $\le 4$, or of order 260--1000 divisible by 4 (except 512 of
rank $\ge 5$ and 768) satisfies (L). Since an example of order $\le 3072$ would contain a case-4 example whose $K$
has order $\le 256$, no example of order $\le 3072$ exists (this covers the orders skipped by the direct search),
and an example, if it exists, has $|G|\ge 6144$.

\subsection{Problem 19.56}\label{app:1956}
Condition (M): $|ab|\ge|a||b|$ for all $\star$-commutators $a,b$ of coprime orders. (M) is inherited by subgroups.
Facts: (b) in a dihedral group $\gen{b,t}$ with $b$ of odd prime order $s$ and $t$ an involution inverting $b$,
$[b,t]=b^{-2}$ generates $\gen b$ and is a $\star$-commutator; (c) in $A_4=V_4\rtimes\gen x$, the commutators
$[x,y]$ ($y\in V_4\setminus 1$) run over $V_4\setminus 1$, so every involution of $A_4$ is a $\star$-commutator;
(d) in $V\rtimes\gen y$ with $V$ elementary abelian and $y$ of prime order acting fixed-point-freely, every
non-trivial element of $V$ is $[v,y]$ for some $v$.

\begin{theorem}\label{thm:1956app}
Every minimal simple group $S$ contains $\star$-commutators $a$ (an involution) and $b$ (of odd prime order $s$) with
$a$ inverting $b$, so $|ab|=2<2s=|a||b|$ and $S$ violates (M). Hence every group with a minimal simple subgroup
violates (M); every counterexample to 19.56 contains a minimal non-soluble subgroup which is again a
counterexample and has non-trivial Frattini subgroup; and a counterexample of minimal order is minimal
non-soluble with non-trivial Frattini subgroup.
\end{theorem}
\begin{proof}
By Thompson~\cite{Thompson}, $S$ is $\PSL(2,2^f)$ ($f\ge 2$ prime), $\PSL(2,3^f)$ ($f$ odd prime), $\PSL(2,p)$ ($p>3$,
$p\equiv\pm2\bmod 5$), $\Sz(2^f)$ ($f$ odd prime) or $\PSL(3,3)$.
For $S=\PSL(2,2^f)$: the Borel subgroup is $V\rtimes C_{q-1}$ with $C_{q-1}$ fixed-point-free on $V\setminus 1$, so by
(d) every involution (one class) is a $\star$-commutator; $S$ has a dihedral subgroup of order $2(q\pm1)$ whose
cyclic part contains an element $b$ of order 3, inverted by an involution $t$; by (b) $b':=b^{-2}$ is a
$\star$-commutator and $a=t$ inverts it. For $S=\PSL(2,q)$, $q$ odd $\ge 5$: $S\supseteq A_4$, so involutions are
$\star$-commutators by (c); some odd prime $s$ divides $(q-1)/2$ or $(q+1)/2$ (both cannot be powers of 2 for $q>3$),
and the dihedral subgroup of order $q\mp1$ provides $b$ of order $s$ and an inverting involution. For
$S=\Sz(q)$: the involutions are the non-trivial elements of $Z(P)$ ($P$ a Sylow 2-subgroup) on which $C_{q-1}$
acts fixed-point-freely, so (d) applies; the normaliser of $C_{q-1}$ is dihedral and $q-1$ has an odd prime factor.
For $\PSL(3,3)$ the pair was found by direct computation (it contains $S_4$ and $D_{12}$). The
remaining assertions follow since (M) passes to subgroups, every non-soluble group contains a minimal non-soluble
subgroup, and a minimal non-soluble group with trivial Frattini subgroup is minimal simple~\cite{Thompson}.
\end{proof}
All 208 non-soluble groups of order $\le 2000$ and all 443 perfect groups of order $\le 10^5$ violate (M)
(\code{work/p19\_56}). The lifting of the configuration through a non-trivial Frattini subgroup is the remaining
obstacle; the agent's partial observations are in the report.

\subsection{Problem 21.26}\label{app:2126}
\begin{proposition}
If $|G|$ has exactly two prime divisors $p,q$, with Sylow subgroups $P$ and $Q$ (so $G=PQ=QP$), there is
$x$ with $P\cap P^x$ and $Q\cap Q^x$ both inclusion-minimal.
\end{proposition}
\begin{proof}
For $a\in P$, $P\cap P^{ga}=(P\cap P^g)^a$, and conjugation by $a$ permutes the family $\{P\cap P^g\}$, so minimality
is invariant under $g\mapsto ga$. Choose $g$ with $P\cap P^g$ minimal and write $g=ba$ ($b\in Q$, $a\in P$); then
$P\cap P^b$ is minimal. Likewise choose $h$ with $Q\cap Q^h$ minimal, $h=a'b'$ ($a'\in P$, $b'\in Q$); then
$Q\cap Q^{a'}$ is minimal. For $x=ba'$: $P\cap P^x=(P\cap P^b)^{a'}$ and $Q\cap Q^x=Q\cap(Q^b)^{a'}=Q\cap Q^{a'}$,
both minimal.
\end{proof}

\subsection{Problem 21.111}\label{app:21111}
\begin{theorem}
For $q=3^n$, $n\ge 2$, $S=\PSL(2,q)$ is a $\tau$-group: for $n$ odd the involutions of $S$, for $n$ even the outer
(diagonal) involutions of $\PGL(2,q)$ are $\tau$-automorphisms.
\end{theorem}
\begin{proof}
For an involution $x\in G=\gen{x,\Inn S}$ and $s\in S$, $\gen{x,x^s}$ is dihedral of order $2\cdot o(xx^s)$, so $x$
is $\tau$ iff no product of two conjugates of $x$ has order divisible by 3. Elements of order divisible by 3 are
unipotent (orders in $S$ are 3 and divisors of $(q\pm1)/2$, and $3\nmid(q\pm1)/2$). If $u=x_1x_2$ with involutions
$x_1,x_2$ conjugate in $G$ then $x_1$ inverts $u$, so $x_1\in N_G(\gen u)$. For $n$ odd, $q\equiv 3\pmod 4$ and
$N_S(\gen u)$ is the Borel subgroup of odd order $q(q-1)/2$, so no involution of $S$ inverts $u$. For $n$ even the
involutions of $\PGL(2,q)$ inverting $u$ lie in its Borel subgroup, whose involutions are conjugate to the image of
$\operatorname{diag}(-1,1)$, which lies in $\PSL(2,q)$ since $-1$ is a square ($q\equiv 1\pmod 4$); so the outer
involutions never invert a unipotent element.
\end{proof}
